\chapter{命令参考}

\section{入口和启动}

直接运行：

\begin{lstlisting}
ccb
\end{lstlisting}

会启动或附着到当前项目 CCB 前台。如果 daemon 未运行，\cmd{ccb} 会自动启动项目 backend，挂载 configured agents，并维护 tmux 前台布局。帮助：

\begin{lstlisting}
ccb --help
ccb ask --help
ccb kill --help
\end{lstlisting}

版本：

\begin{lstlisting}
ccb version
ccb -v
ccb --version
\end{lstlisting}

\section{任务提交}

\begin{lstlisting}
ccb ask [--compact] [--silence] [--chain] [--artifact-request] \
  [--artifact-reply] [--artifact-io] <target> [--] <message...>
\end{lstlisting}

flags：

\begin{itemize}
  \item \cmd{--compact}：请求精简回复。
  \item \cmd{--silence}：成功时静默，失败仍显示。
  \item \cmd{--chain}：把结果作为新任务回调给当前 agent。
  \item \cmd{--artifact-request}：请求体写入 text artifact。
  \item \cmd{--artifact-reply}：回复写入 text artifact。
  \item \cmd{--artifact-io}：同时启用 request 和 reply artifact。
\end{itemize}

示例：

\begin{lstlisting}
ccb ask archi analyze communication flow
ccb ask --compact reviewer review latest diff
ccb ask --silence ops run smoke check
ccb ask --chain archi collect evidence for this task
ccb ask --chain --artifact-reply archi collect long evidence
\end{lstlisting}

active CCB task 中的 nested ask 必须使用 \cmd{--chain} 或 \cmd{--silence}。

\section{ask diagnostics aliases}

\begin{lstlisting}
ccb ask get <job_id>
ccb ask cancel <job_id>
\end{lstlisting}

这些是 diagnostics-oriented aliases。常规取消也可以使用 \cmd{ccb cancel <job_id>}。

\section{观察命令}

\begin{longtable}{p{0.28\linewidth}p{0.62\linewidth}}
\toprule
命令 & 作用 \\
\midrule
\cmd{ccb watch <agent|job_id>} & 轮询 job 或 agent 事件，直到 terminal 或超时。 \\
\cmd{ccb pend <target> [count]} & snapshot 目标状态。 \\
\cmd{ccb pend --watch <target>} & 进入 watch 观察模式。 \\
\cmd{ccb pend --inbox [--detail] <agent>} & 查看 inbox。 \\
\cmd{ccb pend --queue [--detail] <agent>} & 查看 queue。 \\
\cmd{ccb queue [--detail] <target>} & 查看 mailbox queue，target 可为 agent 或 all。 \\
\cmd{ccb inbox [--detail] <agent>} & 查看 agent inbox head 和 pending items。 \\
\cmd{ccb trace <id>} & 从 submission/message/attempt/reply/job id 重建链路。 \\
\bottomrule
\end{longtable}

\section{ack}

\begin{lstlisting}
ccb ack <agent_name> [inbound_event_id]
\end{lstlisting}

ack 只确认 inbox head。它支持 task reply，或者已经 terminal 的 task request。如果指定的 event 不是 head，会被拒绝。已经安排 automatic reply delivery 的 reply 也不能手动 ack。

\section{followup、retry、resubmit、cancel}

\begin{longtable}{p{0.28\linewidth}p{0.62\linewidth}}
\toprule
命令 & 作用 \\
\midrule
\cmd{ccb followup <active_job_id> --message <correction>} & 仅在 provider 能证明 exact active turn 时，把 correction 注入指定 running job；只有 \src{injected} 成功。 \\
\cmd{ccb retry <job_id|attempt_id>} & 在原 message lineage 下追加 retry attempt。要求目标 attempt 已 terminal、不是 completed、且是最新 attempt。 \\
\cmd{ccb resubmit <message_id>} & 从旧 message 创建新的 message/submission/jobs。 \\
\cmd{ccb cancel <job_id>} & 取消 queued 或 active job，并写入 cancelled terminal state。 \\
\bottomrule
\end{longtable}

简单判断：

\begin{itemize}
  \item 修正仍在运行且 provider 支持 exact active-turn injection 的任务，使用 \cmd{followup}；\src{accepted} 表示 durable pending，\src{rejected}、\src{too_late}、\src{terminal} 都不是成功，只有 \src{injected} 成功。
  \item provider/API 临时失败，优先 \cmd{retry}。
  \item 想重新发整条请求，使用 \cmd{resubmit}。
  \item 任务不应继续，使用 \cmd{cancel}。
\end{itemize}

\section{等待命令}

\begin{lstlisting}
ccb wait-any [--timeout N] <target>
ccb wait-all [--timeout N] <target>
ccb wait-quorum [--timeout N] <quorum> <target>
\end{lstlisting}

timeout 必须是正数，quorum 必须是正整数。

\section{生命周期命令}

\begin{longtable}{p{0.30\linewidth}p{0.60\linewidth}}
\toprule
命令 & 作用 \\
\midrule
\cmd{ccb ps} & 查看 agent/backend runtime 状态。 \\
\cmd{ccb ping <agent|all>} & 检查 agent 或 all。 \\
\cmd{ccb clear [agent...]} & 调用 provider-native context clear；交互式 provider 通常接收 \cmd{/clear}，DSH 因没有该命令而轮换 managed native session id。该操作不删除旧 DSH 日志、\src{.ccb} 状态、workspace、auth 或 memory；\cmd{all} 不能和具体 agent 混用。 \\
\cmd{ccb compact [agent...]} & 调用已验证的 native context-compaction 控制（例如 Codex/Claude 的 \cmd{/compact}、Gemini/Qwen/Droid 的 \cmd{/compress}、Crush 的 \cmd{/summarize}）；DSH 通过其结构化 Web API 执行 \cmd{/compact}，不向 host/log pane 输入。busy/queued agent 会被阻止，未验证 provider 返回 \src{unsupported}；不删除任何项目状态。 \\
\cmd{ccb restart <agent>} & 通过 ccbd 重启单个 agent；busy、queued、reply delivery 或 callback continuation 状态会被拒绝；不支持 \cmd{all}。 \\
\cmd{ccb reload [--dry-run]} & 重新加载项目配置；dry-run 只预检。 \\
\cmd{ccb cleanup [--legacy-provider-caches]} & 清理当前停止项目的安全可重建缓存；显式选项额外清理 manifest 校验通过且项目根目录已不存在的旧 Provider 项目缓存。 \\
\cmd{ccb kill [-f|--force]} & 停止 backend 和 configured agents。 \\
\bottomrule
\end{longtable}

\section{诊断和维护命令}

\begin{longtable}{p{0.34\linewidth}p{0.56\linewidth}}
\toprule
命令 & 作用 \\
\midrule
\cmd{ccb doctor} & 常规诊断。 \\
\cmd{ccb doctor --output [PATH]} & 输出诊断 bundle。 \\
\cmd{ccb doctor ps} 或 \cmd{ccb doctor --runtime} & runtime 诊断。 \\
\cmd{ccb doctor logs} 或 \cmd{ccb doctor --logs} & 日志诊断。 \\
\cmd{ccb doctor storage [--json]} & 存储诊断。 \\
\cmd{ccb logs <agent>} & 查看 agent 日志。 \\
\cmd{ccb maintenance [status]} & 查看维护状态。 \\
\cmd{ccb maintenance tick} & 触发一次维护 tick。 \\
\cmd{ccb maintenance tick --force} & 跳过 schedule due 检查，立即运行一次 tick。 \\
\cmd{ccb maintenance tick --no-dispatch} & 只评估和记录，不派发 assessor activation。 \\
\cmd{ccb maintenance schedule --after <duration> [--reason TEXT]} & 设置下一次 heartbeat follow-up 时间。 \\
\cmd{ccb maintenance enable} & 当前实现返回 config-authority 提示；请编辑 \cfg{[maintenance.heartbeat].enabled}。 \\
\cmd{ccb maintenance disable} & 当前实现返回 config-authority 提示；请编辑 \cfg{[maintenance.heartbeat].enabled}。 \\
\bottomrule
\end{longtable}

\section{配置和 fault}

\begin{longtable}{p{0.34\linewidth}p{0.56\linewidth}}
\toprule
命令 & 作用 \\
\midrule
\cmd{ccb config validate} & 验证配置。 \\
\cmd{ccb config approve-commands} & 审阅并批准项目配置中的本地命令字段。 \\
\cmd{ccb fault list} & 列出 fault rules。 \\
\cmd{ccb fault arm <agent> --task-id TASK --reason REASON --count N --error TEXT} & 注入故障。 \\
\cmd{ccb fault clear <rule_id|all>} & 清理故障。 \\
\bottomrule
\end{longtable}

\section{管理和辅助命令}

\begin{longtable}{p{0.34\linewidth}p{0.56\linewidth}}
\toprule
命令 & 作用 \\
\midrule
\cmd{ccb version} & 版本。 \\
\cmd{ccb update} & 更新 CCB，并在交互终端统一检查已安装 provider CLI。 \\
\cmd{ccb update --providers check} & 仅报告 provider 可用更新。 \\
\cmd{ccb update --providers all} & 非交互更新所有可安全管理的已安装 provider。 \\
\cmd{ccb update --providers none} & 本次不检查 provider 更新。 \\
\cmd{ccb update --no-cache-cleanup} & 本次更新跳过旧项目级 Provider 缓存迁移。 \\
\cmd{ccb uninstall} & 卸载。 \\
\cmd{ccb reinstall} & 重装。 \\
\cmd{ccb droid setup-delegation} & 配置 Droid delegation tools。 \\
\cmd{ccb droid test-delegation} & 检查 Droid delegation tools。 \\
\bottomrule
\end{longtable}

交互式 \cmd{ccb update} 只在发现可安全管理的 provider 新版本时提示。
“暂不更新”会在下次执行时再次提示；“跳过此版本”只静默当前准确版本，
更高版本仍会重新提示。Provider 更新不会自动重启正在运行的 pane，新版本
会在下次启动或显式重启后生效。

安装新版本后，新 CCB 会自动清理 manifest 校验通过且项目已删除的旧
Claude/Gemini 项目缓存。当前项目已经停止时会立即清理；活跃项目及其他
现存项目保留到各自下一次成功执行 \cmd{ccb kill}。未知 Provider、损坏
manifest、符号链接、session/auth 和用户级 Gemini 缓存不会被删除。

\section{已移除命令}

\begin{itemize}
  \item \cmd{ccb open} 已移除，直接用 \cmd{ccb}。
  \item \cmd{ccb up} 已移除，agents 由 \src{.ccb/ccb.config} 配置。
  \item \cmd{ccb mail} 和 \cmd{ccb provider} 已移除，使用 \cmd{ask}、\cmd{doctor}、\cmd{trace}。
\end{itemize}
